\subsection{A* Search}
\subsubsection{Idea}
Each state is a partial Futoshiki assignment (\(0\) means unfilled). The solver expands states by
\[
f(s) = g(s) + h(s),
\]
where each action assigns one value to one empty cell (unit cost), \(g(s)\) is the number of assignments already made, and \(h(s)\) estimates the remaining assignments.

To reduce branching, the next cell is selected by MRV (smallest domain). Candidate values are filtered by row/column uniqueness and inequality constraints; optional AC-3 further removes unsupported values. This combines informed search (A*) with CSP-style pruning.

\subsubsection{Heuristic Function and Admissibility Proof}
Let \(U(s)\) be the number of unassigned cells in state \(s\). We use:
\[
h(s)=U(s)=N^2-A(s),
\]
where \(A(s)\) is the number of assigned cells.

AC-3 is used as a separate feasibility check: if any domain becomes empty, the state is discarded.

\textbf{Admissibility:}
\begin{itemize}
\item Any complete solution from state \(s\) must assign all \(U(s)\) empty cells.
\item Each action can assign at most one cell.
\item Hence the true remaining cost \(h^*(s)\) satisfies \(h^*(s) \ge U(s)=h(s)\).
\end{itemize}
So \(h(s)\) never overestimates and is admissible.

\textbf{Consistency:}
\begin{itemize}
\item For one valid move \(s\rightarrow s'\), \(U(s')=U(s)-1\).
\item With unit step cost \(c(s,s')=1\):
\[
h(s)=U(s)=1+U(s')=c(s,s')+h(s').
\]
\item Therefore \(h\) is consistent.
\end{itemize}

\subsubsection{Pseudocode}
\begin{algorithm}[H]
\caption{A* for Futoshiki (with MRV and AC-3)}
\begin{algorithmic}[1]
\State Build initial state \(s_0\) from puzzle
\State Validate givens; if invalid then return \textbf{failure}
\State Compute domains and heuristic for \(s_0\)
\State Initialize \textit{closed} as an empty set
\State Initialize \textit{best\_g}[\(s_0\)]
\State Push \(s_0\) into priority queue by key \((f, h, -g)\)
\While{priority queue not empty}
\State Pop best state \(s\)
\If{\(s \in\) \textit{closed}} \State \textbf{continue} \EndIf
\State Add \(s\) to \textit{closed}
\If{\(s\) is goal} \Return solved grid \EndIf
\State Select unassigned cell \((r,c)\) by MRV
\For{each value \(v\) in domain\((r,c)\)}
\If{assigning \(v\) is consistent}
\State Create child state \(s'\)
\If{\(s' \in\) \textit{closed}} \State \textbf{continue} \EndIf
\State Compute \(g(s')\)
\If{\textit{best\_g}[\(s'\)] exists and \(g(s')\) is not better}
\State \textbf{continue}
\EndIf
\State Recompute domains (and AC-3 if enabled)
\If{no empty domain in \(s'\)}
\State Set \(g(s') = g(s)+1\), evaluate \(h(s')\), \(f(s')=g(s')+h(s')\)
\State Update \textit{best\_g}[\(s'\)]
\State Push \(s'\) into priority queue
\EndIf
\EndIf
\EndFor
\EndWhile
\State \Return \textbf{failure}
\end{algorithmic}
\end{algorithm}

\subsubsection{Illustration}
For a \(4 \times 4\) instance, the solver:
\begin{enumerate}
\item Initializes domains for empty cells.
\item Applies consistency filtering (and AC-3 if enabled).
\item Chooses the most constrained cell by MRV.
\item Expands assignments by priority \(f=g+h\).
\end{enumerate}

Example progression:
\begin{itemize}
\item State \(s_0\): clue cells fixed, inequalities loaded.
\item State \(s_1\): assign MRV cell, domains shrink.
\item State \(s_2\): additional forced assignment appears.
\item Goal: all cells assigned; row/column uniqueness and inequalities satisfied.
\end{itemize}

\subsubsection{Algorithm Analysis}
Let \(N\) be board size (\(N^2\) cells) and \(b\) the average branching factor after filtering.

\textbf{Time complexity (worst case):}
\begin{itemize}
\item Exponential search in remaining depth, approximately \(O(b^d)\) with \(d \le N^2\).
\item AC-3 adds polynomial overhead per state; runtime is mainly determined by expanded states.
\end{itemize}

\textbf{Space complexity:}
\begin{itemize}
\item A* stores open/closed states and domains, so memory use is higher than depth-first methods.
\item Worst-case space is also exponential.
\end{itemize}

\textbf{Practical behavior:}
\begin{itemize}
\item \(h(s)=U(s)\) is admissible and simple to justify.
\item Most speedup comes from MRV and AC-3 pruning.
\item \textit{closed} avoids re-expansion; \textit{best\_g} removes inferior duplicate paths.
\item Works well on medium instances, but memory pressure can rise on larger or weakly constrained puzzles.
\end{itemize}
